% !TEX program = pdflatex
% ------------------------------------------------------------
% MATH 431 Homework Template (plug-and-chug friendly)
%
% Goals:
% - Same clean header style as your MATH 521 template.
% - Easy “Problem -> Solution” workflow.
% - Lightweight (no tcolorbox); uses mdframed for optional boxes.
% - Good for probability homework: sample space, probability measure, events.
%
% Quick use:
%   1) Edit the Metadata block.
%   2) For each exercise, use \problem{<section>.<number>} and paste the prompt.
%   3) Write your solution in the solution environment.
% ------------------------------------------------------------

\documentclass[12pt]{article}

% =========================
% 0) Metadata (EDIT ME)
% =========================
\newcommand{\CourseCode}{MATH 431}
\newcommand{\CourseTitle}{Introduction to Probability}
\newcommand{\CourseSection}{LEC 002} % change to 001 or 003 if needed
\newcommand{\Semester}{Spring 2026}
\newcommand{\Instructor}{Sarah Tammen} % LEC 001: Ander Aguirre | LEC 002: Sarah Tammen | LEC 003: Nicholas Derr
\newcommand{\MeetingInfo}{MWF 9:55--10:45, 6102 Sewell Social Sciences} % LEC 001 default; update if in 002/003
\newcommand{\HomeworkNumber}{9}
\newcommand{\YourName}{Alejandro De La Torre}
\newcommand{\DueDate}{March 27, 2026} % or e.g. February 2, 2026

% =========================
% 1) Core math tools
% =========================
\usepackage{amsmath, amssymb, amsthm}

% =========================
% 2) Lists and spacing
% =========================
\usepackage{enumitem}
\setlist[enumerate,1]{label=\textbf{(\alph*)}, leftmargin=2em, itemsep=0.25em}
\setlist[itemize]{leftmargin=1.5em, itemsep=0.25em}
\setlength{\parskip}{0.35em}
\setlength{\parindent}{0pt}

% =========================
% 3) Page geometry + header/footer
% =========================
\usepackage{geometry}
\geometry{margin=1in}

\usepackage{fancyhdr}
\pagestyle{fancy}
\fancyhf{}
\lhead{\CourseCode\ --- \CourseSection, \Semester}
\rhead{Homework \HomeworkNumber}
\cfoot{\thepage}
\renewcommand{\headrulewidth}{0.4pt}
\setlength{\headheight}{15pt}
\addtolength{\topmargin}{-2pt}

% =========================
% 4) Links (subtle)
% =========================
\usepackage[hidelinks]{hyperref}
\setcounter{tocdepth}{1}

% =========================
% 5) Figures (optional)
% =========================
\usepackage{graphicx}
\graphicspath{{images/}}

% =========================
% 6) Lightweight boxes (optional)
% =========================
\usepackage{mdframed}
\newmdenv[
  skipabove=0.6\baselineskip,
  skipbelow=0.6\baselineskip,
  linewidth=0.6pt,
  roundcorner=4pt,
  innerleftmargin=10pt,
  innerrightmargin=10pt,
  innertopmargin=8pt,
  innerbottommargin=8pt
]{boxnote}

% =========================
% 7) Probability / notation macros
% =========================
\newcommand{\R}{\mathbb{R}}
\newcommand{\N}{\mathbb{N}}
\newcommand{\Z}{\mathbb{Z}}

\newcommand{\OmegaSpace}{\Omega}
\newcommand{\FField}{\mathcal{F}}
% Probability measure in case it is relevant or want to distinguish it from P(A)
%\newcommand{\PMeasure}{\mathbb{P}} 
\newcommand{\E}{\mathbb{E}}

\newcommand{\given}{\,\middle|\,}
\newcommand{\ind}{\mathbf{1}}

% Common “defined as” symbol
\newcommand{\defeq}{\vcentcolon=}

% =========================
% 8) Problem command (TOC + PDF bookmarks)
% Usage:
%   \problem{<label>}
%   \problem[Short title]{<label>}
% Example:
%   \problem{1.4}
%   \problem[Sample space]{1.4}
% =========================
\makeatletter
\newcommand{\problem}[2][]{%
  \def\prob@title{Problem #2\if\relax\detokenize{#1}\relax\else\ (\textit{#1})\fi}%
  \phantomsection%
  \pdfbookmark[1]{\prob@title}{prob#2}%
  \addcontentsline{toc}{section}{\prob@title}%
  \section*{\prob@title}%
  \vspace{-0.25em}\hrule\vspace{0.8em}%
}
\makeatother

% =========================
% 9) Solution environment
% =========================
\newenvironment{solution}{%
  \noindent\textbf{Solution.}\ \ignorespaces
}{\par}

% Optional scratch / planning block
\newenvironment{scratch}{%
  \begin{boxnote}\textbf{Scratch / Setup.}\ \small
}{\end{boxnote}}

% Final answer command: expects math only (no $$ or \[ \])
\newcommand{\finalanswer}[1]{%
  \par\medskip
  \noindent\textbf{Final:}\ \fbox{$#1$}\par
} % AGAIN: MATH ONLY INSIDE THE BOX; if word answer, use \text{} inside or use \fbox{TEXT} or \framebox

% =========================
% 10) Title block
% =========================
\title{Homework \HomeworkNumber}
\author{\YourName \\
\CourseCode\ --- \CourseTitle \\
\CourseSection \\
Instructor: \Instructor}
\date{Due: \DueDate}

\begin{document}
\maketitle

% \section*{Collaboration / Resources}
% (Optional) Write “I worked alone” or list collaborators/resources.

\tableofcontents
\bigskip
\hrule
\bigskip
\newpage

% ============================================================
% Problems (duplicate blocks as needed)
% ============================================================


\problem{6.2}
The joint probability mass function of the random variables $(X,Y)$ is given by the following table:
\[
\begin{array}{c|cccc}
 & \multicolumn{4}{c}{Y} \\
X & 0 & 1 & 2 & 3 \\ \hline
1 & \frac{1}{15} & \frac{1}{15} & \frac{2}{15} & \frac{1}{15} \\
2 & \frac{1}{10} & \frac{1}{10} & \frac{1}{5} & \frac{1}{10} \\
3 & \frac{1}{30} & \frac{1}{30} & 0 & \frac{1}{10}
\end{array}
\]
\begin{enumerate}[label=(\alph*)]
\item Find the marginal probability mass functions of $X$ and $Y$.
\item Calculate the probability $P(X+Y^2\le 2)$.
\end{enumerate}

\begin{solution}
\begin{scratch}
Let $p_{X,Y}(x,y)$ denote the table entries. For part (a),
\[
p_X(x)=\sum_y p_{X,Y}(x,y),\qquad p_Y(y)=\sum_x p_{X,Y}(x,y).
\]
For part (b), we list the pairs with $x+y^2\le 2$. Since $y^2\ge 4$ when
$y=2$ or $3$, only $y=0,1$ can contribute.
\end{scratch}

\begin{enumerate}[label=(\alph*)]
\item Summing the rows gives
\[
p_X(1)=\frac{1}{15}+\frac{1}{15}+\frac{2}{15}+\frac{1}{15}=\frac{1}{3},
\]
\[
p_X(2)=\frac{1}{10}+\frac{1}{10}+\frac{1}{5}+\frac{1}{10}=\frac{1}{2},
\]
\[
p_X(3)=\frac{1}{30}+\frac{1}{30}+0+\frac{1}{10}=\frac{1}{6}.
\]
So the marginal p.m.f. of $X$ is
\[
\begin{array}{c|ccc}
x & 1 & 2 & 3 \\ \hline
p_X(x) & \frac13 & \frac12 & \frac16
\end{array}
\]
Summing the columns gives
\[
p_Y(0)=\frac{1}{15}+\frac{1}{10}+\frac{1}{30}=\frac{1}{5},
\qquad
p_Y(1)=\frac{1}{15}+\frac{1}{10}+\frac{1}{30}=\frac{1}{5},
\]
\[
p_Y(2)=\frac{2}{15}+\frac{1}{5}+0=\frac{1}{3},
\qquad
p_Y(3)=\frac{1}{15}+\frac{1}{10}+\frac{1}{10}=\frac{4}{15}.
\]
Hence
\[
\begin{array}{c|cccc}
y & 0 & 1 & 2 & 3 \\ \hline
p_Y(y) & \frac15 & \frac15 & \frac13 & \frac{4}{15}
\end{array}
\]
\item The condition $x+y^2\le 2$ is satisfied only for
\[
(x,y)=(1,0),\ (1,1),\ (2,0).
\]
Therefore,
\[
P(X+Y^2\le 2)=p_{X,Y}(1,0)+p_{X,Y}(1,1)+p_{X,Y}(2,0)
=\frac{1}{15}+\frac{1}{15}+\frac{1}{10}
=\frac{7}{30}.
\]
\end{enumerate}

\finalanswer{\text{(a)}\ p_X(1)=\frac13,\ p_X(2)=\frac12,\ p_X(3)=\frac16}
\finalanswer{\text{(a)}\ p_Y(0)=\frac15,\ p_Y(1)=\frac15,\ p_Y(2)=\frac13,\ p_Y(3)=\frac{4}{15}}
\finalanswer{\text{(b)}\ P(X+Y^2\le 2)=\frac{7}{30}}

\end{solution}

\newpage

\problem{6.6}
Suppose that $X,Y$ are jointly continuous with joint probability density function
\[
f(x,y)=
\begin{cases}
xe^{-x(1+y)}, & \text{if } x>0 \text{ and } y>0 \\
0, & \text{otherwise.}
\end{cases}
\]
\begin{enumerate}[label=(\alph*)]
\item Find the marginal density functions of $X$ and $Y$.
\item Calculate the expectation $E[XY]$.
\item Calculate the expectation $E\!\left[\frac{X}{1+Y}\right]$.
\end{enumerate}

\begin{solution}
\begin{scratch}
The support is $x>0$, $y>0$. For the marginals,
\[
f_X(x)=\int_0^\infty x e^{-x(1+y)}\,dy,\qquad
f_Y(y)=\int_0^\infty x e^{-x(1+y)}\,dx.
\]
Useful integrals:
\[
\int_0^\infty e^{-ay}\,dy=\frac1a,\quad
\int_0^\infty xe^{-ax}\,dx=\frac1{a^2},\quad
\int_0^\infty x^2e^{-ax}\,dx=\frac{2}{a^3}
\qquad(a>0).
\]
\end{scratch}

\begin{enumerate}[label=(\alph*)]
\item For $x>0$,
\[
f_X(x)=\int_0^\infty x e^{-x(1+y)}\,dy
=xe^{-x}\int_0^\infty e^{-xy}\,dy
=xe^{-x}\cdot \frac1x
=e^{-x}.
\]
Thus
\[
f_X(x)=
\begin{cases}
e^{-x}, & x>0,\\
0, & \text{otherwise.}
\end{cases}
\]
For $y>0$,
\[
f_Y(y)=\int_0^\infty x e^{-x(1+y)}\,dx
=\frac{1}{(1+y)^2}.
\]
Hence
\[
f_Y(y)=
\begin{cases}
\dfrac{1}{(1+y)^2}, & y>0,\\
0, & \text{otherwise.}
\end{cases}
\]
\item
\[
E[XY]=\int_0^\infty\int_0^\infty xy \, x e^{-x(1+y)}\,dy\,dx
=\int_0^\infty x^2e^{-x}\left(\int_0^\infty y e^{-xy}\,dy\right)dx.
\]
Since $\int_0^\infty y e^{-xy}\,dy=\frac{1}{x^2}$ for $x>0$,
\[
E[XY]=\int_0^\infty e^{-x}\,dx=1.
\]
\item
\[
E\!\left[\frac{X}{1+Y}\right]
=\int_0^\infty\int_0^\infty \frac{x}{1+y}\,x e^{-x(1+y)}\,dy\,dx.
\]
Integrating first in $x$,
\[
E\!\left[\frac{X}{1+Y}\right]
=\int_0^\infty \frac{1}{1+y}\left(\int_0^\infty x^2 e^{-x(1+y)}\,dx\right)dy
\]
\[
=\int_0^\infty \frac{1}{1+y}\cdot \frac{2}{(1+y)^3}\,dy
=2\int_0^\infty \frac{1}{(1+y)^4}\,dy
=\frac{2}{3}.
\]
\end{enumerate}

\finalanswer{\text{(a)}\ f_X(x)=e^{-x}\ind_{\{x>0\}},\quad f_Y(y)=\frac{1}{(1+y)^2}\ind_{\{y>0\}}}
\finalanswer{\text{(b)}\ E[XY]=1,\qquad \text{(c)}\ E\!\left[\frac{X}{1+Y}\right]=\frac{2}{3}}

\end{solution}

\newpage

\problem{6.8}
Determine whether the following pairs of random variables $X$ and $Y$ are independent.
\begin{enumerate}[label=(\alph*)]
\item Random variables $X$ and $Y$ of Exercise 6.2.
\item Random variables $X$ and $Y$ of Exercise 6.5.
\setcounter{enumi}{3}
\item Random variables $X$ and $Y$ of Exercise 6.7.
\end{enumerate}

\begin{solution}
\begin{scratch}
For independence we compare the joint distribution with the product of the
marginals.
\begin{itemize}
\item In part (a), it is enough to find one pair $(x,y)$ with
$p_{X,Y}(x,y)\ne p_X(x)p_Y(y)$.
\item In part (b), use the marginals from Exercise 6.5:
$f_X(x)=\frac67x+\frac47$ and $f_Y(y)=\frac{12}{7}y^2+\frac67y$ on $[0,1]$.
\item In part (d), the support is the triangle
$\{(x,y):0\le x,\ 0\le y,\ x+y\le 1\}$, not a rectangle.
\end{itemize}
\end{scratch}

\begin{enumerate}[label=(\alph*)]
\item From Exercise 6.2,
\[
P(X=3)=\frac16>0,\qquad P(Y=2)=\frac13>0,
\]
but from the table,
\[
P(X=3,Y=2)=0.
\]
Hence $P(X=3,Y=2)\ne P(X=3)P(Y=2)$, so $X$ and $Y$ are not independent.
\item From Exercise 6.5,
\[
f_{X,Y}(x,y)=\frac{12}{7}(xy+y^2),\qquad 0\le x\le 1,\ 0\le y\le 1,
\]
and the marginals are
\[
f_X(x)=\frac67x+\frac47,\qquad f_Y(y)=\frac{12}{7}y^2+\frac67y.
\]
Evaluating at $x=y=\frac14$,
\[
f_{X,Y}\!\left(\frac14,\frac14\right)=\frac{12}{7}\left(\frac1{16}+\frac1{16}\right)=\frac{3}{14},
\]
while
\[
f_X\!\left(\frac14\right)f_Y\!\left(\frac14\right)
=\left(\frac{11}{14}\right)\left(\frac{9}{28}\right)
=\frac{99}{392}\ne \frac{3}{14}.
\]
So $X$ and $Y$ are not independent.
\item In Exercise 6.7, $(X,Y)$ is uniform on the triangle
\[
\{(x,y):0\le x,\ 0\le y,\ x+y\le 1\}.
\]
The marginal densities are positive on $(0,1)$:
\[
f_X(x)=2(1-x),\qquad f_Y(y)=2(1-y).
\]
Take $x=y=\frac34$. Then
\[
f_X\!\left(\frac34\right)>0,\qquad f_Y\!\left(\frac34\right)>0,
\]
but
\[
f_{X,Y}\!\left(\frac34,\frac34\right)=0
\]
because $\frac34+\frac34>1$, so the point lies outside the triangle. Therefore
$f_{X,Y}(x,y)\ne f_X(x)f_Y(y)$, and $X$ and $Y$ are not independent.
\end{enumerate}

\finalanswer{\text{(a), (b), and (d) are all not independent}}

\end{solution}

\newpage

\problem{6.18}
Suppose that $X$ and $Y$ are integer-valued random variables with joint probability mass function given by
\[
p_{X,Y}(a,b)=
\begin{cases}
\dfrac{1}{4a}, & \text{for } 1\le b\le a\le 4 \\
0, & \text{otherwise.}
\end{cases}
\]
\begin{enumerate}[label=(\alph*)]
\item Show that this is indeed a joint probability mass function.
\item Find the marginal probability mass functions of $X$ and $Y$.
\item Find $P(X=Y+1)$.
\end{enumerate}

\begin{solution}
\begin{scratch}
The support is
\[
\{(a,b)\in \mathbb{Z}^2:1\le b\le a\le 4\}.
\]
For part (a),
\[
\sum_{a=1}^4\sum_{b=1}^a \frac{1}{4a}
=\sum_{a=1}^4 \frac{a}{4a}=1.
\]
For the marginals,
\[
p_X(a)=\sum_{b=1}^a p_{X,Y}(a,b),\qquad
p_Y(b)=\sum_{a=b}^4 p_{X,Y}(a,b).
\]
\end{scratch}

\begin{enumerate}[label=(\alph*)]
\item Since $p_{X,Y}(a,b)\ge 0$ on its support, it remains only to check that the
total mass is $1$:
\[
\sum_{a=1}^4\sum_{b=1}^a \frac{1}{4a}
=\sum_{a=1}^4 \frac{a}{4a}
=4\cdot \frac14
=1.
\]
So this is a valid joint probability mass function.
\item For $a=1,2,3,4$,
\[
p_X(a)=\sum_{b=1}^a \frac{1}{4a}=a\cdot \frac{1}{4a}=\frac14.
\]
Thus
\[
p_X(a)=\frac14,\qquad a=1,2,3,4.
\]
For $Y$,
\[
p_Y(1)=\frac14+\frac18+\frac1{12}+\frac1{16}=\frac{25}{48},
\]
\[
p_Y(2)=\frac18+\frac1{12}+\frac1{16}=\frac{13}{48},
\]
\[
p_Y(3)=\frac1{12}+\frac1{16}=\frac{7}{48},
\qquad
p_Y(4)=\frac1{16}.
\]
\item The event $\{X=Y+1\}$ consists of the admissible pairs
\[
(2,1),\ (3,2),\ (4,3).
\]
Therefore,
\[
P(X=Y+1)=\frac18+\frac1{12}+\frac1{16}
=\frac{6+4+3}{48}
=\frac{13}{48}.
\]
\end{enumerate}

\finalanswer{\text{(a) valid joint pmf},\quad p_X(a)=\frac14\ \text{for } a=1,2,3,4}
\finalanswer{p_Y(1)=\frac{25}{48},\ p_Y(2)=\frac{13}{48},\ p_Y(3)=\frac{7}{48},\ p_Y(4)=\frac{1}{16}}
\finalanswer{\text{(c)}\ P(X=Y+1)=\frac{13}{48}}

\end{solution}

\newpage

\problem{6.24}
An urn contains 1 green ball, 1 red ball, 1 yellow ball and 1 white ball. I draw 3 balls with replacement. What is the probability that exactly two balls are of the same color?

\begin{solution}
\begin{scratch}
Treat each draw as a color outcome from $\{G,R,Y,W\}$. There are
\[
4^3=64
\]
equally likely ordered color sequences. For a favorable outcome, one color
appears exactly twice and a different color appears once.
\end{scratch}

To get exactly two balls of one color and one ball of a different color:
\[
\text{choose the repeated color: }4\text{ ways,}
\]
\[
\text{choose the singleton color: }3\text{ ways,}
\]
\[
\text{choose the position of the singleton: }3\text{ ways.}
\]
Hence the number of favorable ordered outcomes is
\[
4\cdot 3\cdot 3=36.
\]
Since there are $4^3=64$ possible ordered outcomes in total,
\[
P(\text{exactly two balls are of the same color})=\frac{36}{64}=\frac{9}{16}.
\]

\finalanswer{\frac{9}{16}}

\end{solution}

\newpage

\problem{6.27}
Suppose that $X_1$ and $X_2$ are independent random variables with
\[
P(X_1=1)=P(X_1=-1)=\frac{1}{2}
\]
and
\[
P(X_2=1)=1-P(X_2=-1)=p
\]
with $0<p<1$. Let $Y=X_1X_2$. Show that $X_2$ and $Y$ are independent.

\begin{solution}
\begin{scratch}
Both $X_2$ and $Y=X_1X_2$ take values in $\{-1,1\}$. To prove independence,
it is enough to compute
\[
P(X_2=a,Y=b)
\]
for $a,b\in\{-1,1\}$ and show it factors as
$P(X_2=a)P(Y=b)$.
\end{scratch}

Fix $a,b\in\{-1,1\}$. Since $Y=X_1X_2$, the event $\{X_2=a,Y=b\}$ is the
same as
\[
\{X_2=a,\ X_1X_2=b\}
=\{X_2=a,\ X_1=b/a\}.
\]
Therefore, by independence of $X_1$ and $X_2$,
\[
P(X_2=a,Y=b)=P(X_2=a)P(X_1=b/a).
\]
Because $b/a\in\{-1,1\}$ and
\[
P(X_1=1)=P(X_1=-1)=\frac12,
\]
we get
\[
P(X_2=a,Y=b)=P(X_2=a)\cdot \frac12.
\]
It remains to compute $P(Y=b)$. We have
\[
P(Y=1)=P(X_1=X_2)
=P(X_1=1,X_2=1)+P(X_1=-1,X_2=-1)
\]
\[
=\frac12 p+\frac12(1-p)=\frac12.
\]
Similarly,
\[
P(Y=-1)=1-P(Y=1)=\frac12.
\]
Hence
\[
P(X_2=a,Y=b)=P(X_2=a)\cdot \frac12
=P(X_2=a)P(Y=b)
\]
for all $a,b\in\{-1,1\}$. Therefore $X_2$ and $Y$ are independent.

\finalanswer{X_2\ \text{and}\ Y\ \text{are independent}}

\end{solution}

\newpage

\problem{6.28}
Let $X$ and $Y$ be independent Geom$(p)$ random variables. Let
\[
V=\min(X,Y)
\]
and
\[
W=
\begin{cases}
0, & \text{if } X<Y \\
1, & \text{if } X=Y \\
2, & \text{if } X>Y.
\end{cases}
\]
Find the joint probability mass function of $V$ and $W$ and show that $V$ and $W$ are independent.

\emph{Hint.} Use the joint probability mass function of $X$ and $Y$ to compute the joint probability mass function of $V$ and $W$.

\begin{solution}
\begin{scratch}
Let $q=1-p$. Since $X,Y\sim\mathrm{Geom}(p)$ independently, for $k\ge 1$,
\[
P(X=k)=P(Y=k)=pq^{k-1}.
\]
For $v\ge 1$:
\[
P(V=v,W=0)=P(X=v,\ Y>v),
\]
\[
P(V=v,W=1)=P(X=v,\ Y=v),
\]
\[
P(V=v,W=2)=P(Y=v,\ X>v).
\]
\end{scratch}

For $v=1,2,\dots$,
\[
P(V=v,W=0)=P(X=v)P(Y>v)=pq^{v-1}\cdot q^v=pq^{2v-1},
\]
using independence and $P(Y>v)=q^v$.
Similarly,
\[
P(V=v,W=2)=pq^{2v-1}.
\]
Also,
\[
P(V=v,W=1)=P(X=v,Y=v)=p^2q^{2v-2}.
\]
Thus the joint p.m.f. is
\[
p_{V,W}(v,w)=
\begin{cases}
pq^{2v-1}, & v\ge 1,\ w=0,\\
p^2q^{2v-2}, & v\ge 1,\ w=1,\\
pq^{2v-1}, & v\ge 1,\ w=2,\\
0, & \text{otherwise.}
\end{cases}
\]

Now compute the marginals. For $v\ge 1$,
\[
p_V(v)=pq^{2v-1}+p^2q^{2v-2}+pq^{2v-1}
=q^{2v-2}\bigl(2pq+p^2\bigr)
=p(2-p)q^{2v-2}.
\]
For $W$,
\[
P(W=1)=\sum_{v=1}^\infty p^2q^{2v-2}
=p^2\sum_{j=0}^\infty q^{2j}
=\frac{p^2}{1-q^2}
=\frac{p}{2-p}.
\]
By symmetry, $P(W=0)=P(W=2)$, and so
\[
P(W=0)=P(W=2)=\frac{1-P(W=1)}{2}
=\frac{1-p}{2-p}.
\]

Finally, for $v\ge 1$,
\[
p_V(v)P(W=0)=p(2-p)q^{2v-2}\cdot \frac{1-p}{2-p}=pq^{2v-1}=p_{V,W}(v,0),
\]
\[
p_V(v)P(W=1)=p(2-p)q^{2v-2}\cdot \frac{p}{2-p}=p^2q^{2v-2}=p_{V,W}(v,1),
\]
and similarly $p_V(v)P(W=2)=p_{V,W}(v,2)$. Hence $V$ and $W$ are independent.

\finalanswer{p_{V,W}(v,w)=
\begin{cases}
pq^{2v-1}, & v\ge 1,\ w=0,\\
p^2q^{2v-2}, & v\ge 1,\ w=1,\\
pq^{2v-1}, & v\ge 1,\ w=2,\\
0, & \text{otherwise}
\end{cases}}
\finalanswer{p_V(v)=p(2-p)q^{2v-2}\ \text{for } v\ge 1}
\finalanswer{P(W=0)=P(W=2)=\frac{1-p}{2-p},\quad P(W=1)=\frac{p}{2-p}}
\finalanswer{V\ \text{and}\ W\ \text{are independent}}

\end{solution}

\newpage

\problem{6.30}
Suppose that $X\sim \text{Geom}(p)$ and $Y\sim \text{Poisson}(\lambda)$ are independent. Find the probability $P(X-1=Y)$.

\begin{solution}
\begin{scratch}
Since $Y\sim\mathrm{Poisson}(\lambda)$ takes values $0,1,2,\dots$, the event
$\{X-1=Y\}$ is the disjoint union
\[
\bigcup_{k=0}^\infty \{X=k+1,\ Y=k\}.
\]
Use independence and the p.m.f.s
\[
P(X=k+1)=p(1-p)^k,\qquad P(Y=k)=e^{-\lambda}\frac{\lambda^k}{k!}.
\]
\end{scratch}

By independence,
\[
P(X-1=Y)=\sum_{k=0}^\infty P(X=k+1)P(Y=k)
\]
\[
=\sum_{k=0}^\infty p(1-p)^k e^{-\lambda}\frac{\lambda^k}{k!}
=pe^{-\lambda}\sum_{k=0}^\infty \frac{(\lambda(1-p))^k}{k!}.
\]
Recognizing the exponential series,
\[
P(X-1=Y)=pe^{-\lambda}e^{\lambda(1-p)}=pe^{-\lambda p}.
\]

\finalanswer{P(X-1=Y)=pe^{-\lambda p}}

\end{solution}

\newpage

\problem{6.34}
Let $(X,Y)$ be a uniformly distributed random point on the quadrilateral $D$ with vertices $(0,0)$, $(2,0)$, $(1,1)$ and $(0,1)$.
\begin{enumerate}[label=(\alph*)]
\item Find the joint density function of $(X,Y)$ and the marginal density functions of $X$ and $Y$.
\item Find $E[X]$ and $E[Y]$.
\item Are $X$ and $Y$ independent?
\end{enumerate}

\begin{solution}
\begin{scratch}
The quadrilateral can be described as
\[
D=\{(x,y):0\le y\le 1,\ 0\le x\le 2-y\}.
\]
Its area is
\[
\int_0^1 (2-y)\,dy=\frac32.
\]
So the joint density is constant $2/3$ on $D$.
For marginals:
\begin{itemize}
\item If $0\le x\le 1$, then $y$ ranges from $0$ to $1$.
\item If $1\le x\le 2$, then $y$ ranges from $0$ to $2-x$.
\item If $0\le y\le 1$, then $x$ ranges from $0$ to $2-y$.
\end{itemize}
\end{scratch}

\begin{enumerate}[label=(\alph*)]
\item Since $(X,Y)$ is uniform on $D$,
\[
f_{X,Y}(x,y)=
\begin{cases}
\dfrac23, & (x,y)\in D,\\
0, & \text{otherwise.}
\end{cases}
\]
Equivalently,
\[
f_{X,Y}(x,y)=
\begin{cases}
\dfrac23, & 0\le y\le 1,\ 0\le x\le 2-y,\\
0, & \text{otherwise.}
\end{cases}
\]
For the marginal of $X$,
\[
f_X(x)=
\begin{cases}
\int_0^1 \frac23\,dy=\frac23, & 0\le x\le 1,\\[8pt]
\int_0^{2-x}\frac23\,dy=\frac43-\frac23x, & 1\le x\le 2,\\[8pt]
0, & \text{otherwise.}
\end{cases}
\]
For the marginal of $Y$,
\[
f_Y(y)=
\begin{cases}
\int_0^{2-y}\frac23\,dx=\frac43-\frac23y, & 0\le y\le 1,\\[8pt]
0, & \text{otherwise.}
\end{cases}
\]
\item Using the marginals,
\[
E[X]=\int_0^1 x\cdot \frac23\,dx+\int_1^2 x\left(\frac43-\frac23x\right)dx
=\frac{7}{9}.
\]
Also,
\[
E[Y]=\int_0^1 y\left(\frac43-\frac23y\right)dy=\frac49.
\]
\item $X$ and $Y$ are not independent. For instance,
\[
P\!\left(X>\frac32,\ Y>\frac12\right)=0
\]
because every point in $D$ satisfies $x+y\le 2$, so $x>\frac32$ and
$y>\frac12$ cannot happen together. But
\[
P\!\left(X>\frac32\right)>0,\qquad P\!\left(Y>\frac12\right)>0,
\]
so the intersection probability is not the product of the marginals.
\end{enumerate}

\finalanswer{\text{(a)}\ f_{X,Y}(x,y)=\frac23\ \text{for } (x,y)\in D,\ 0\ \text{otherwise}}
\finalanswer{f_X(x)=\begin{cases}\frac23,&0\le x\le 1\\ \frac43-\frac23x,&1\le x\le 2\\0,&\text{otherwise}\end{cases}}
\finalanswer{f_Y(y)=\begin{cases}\frac43-\frac23y,&0\le y\le 1\\0,&\text{otherwise}\end{cases}}
\finalanswer{\text{(b)}\ E[X]=\frac79,\ E[Y]=\frac49,\qquad \text{(c)}\ X\ \text{and}\ Y\ \text{are not independent}}

\end{solution}

\newpage

\problem{6.36}
Suppose that $X,Y$ are jointly continuous with joint probability density function
\[
f(x,y)=ce^{-x^2/2-(x-y)^2/2},\qquad x,y\in(-\infty,\infty),
\]
for some constant $c$.
\begin{enumerate}[label=(\alph*)]
\item Find the value of the constant $c$.
\item Find the marginal density functions of $X$ and $Y$.
\item Determine whether $X$ and $Y$ are independent.
\end{enumerate}
% Homework hint: use that the integral of the normal density is equal to 1.

\begin{solution}
\begin{scratch}
For normalization,
\[
1=\iint_{\mathbb{R}^2} c e^{-x^2/2-(x-y)^2/2}\,dy\,dx.
\]
Integrate first in $y$ and use the hint:
\[
\int_{-\infty}^{\infty} e^{-(x-y)^2/2}\,dy=\sqrt{2\pi}.
\]
For $f_Y$, complete the square:
\[
\frac{x^2}{2}+\frac{(x-y)^2}{2}
=x^2-xy+\frac{y^2}{2}
=\left(x-\frac y2\right)^2+\frac{y^2}{4}.
\]
\end{scratch}

\begin{enumerate}[label=(\alph*)]
\item We compute
\[
1=c\int_{-\infty}^{\infty}\int_{-\infty}^{\infty}
e^{-x^2/2-(x-y)^2/2}\,dy\,dx.
\]
For fixed $x$,
\[
\int_{-\infty}^{\infty} e^{-(x-y)^2/2}\,dy=\sqrt{2\pi},
\]
so
\[
1=c\sqrt{2\pi}\int_{-\infty}^{\infty} e^{-x^2/2}\,dx
=c\sqrt{2\pi}\sqrt{2\pi}=c(2\pi).
\]
Hence
\[
c=\frac{1}{2\pi}.
\]
\item Using the value of $c$,
\[
f_X(x)=\int_{-\infty}^{\infty}\frac{1}{2\pi}
e^{-x^2/2-(x-y)^2/2}\,dy
=\frac{1}{2\pi}e^{-x^2/2}\int_{-\infty}^{\infty}e^{-(x-y)^2/2}\,dy
\]
\[
=\frac{1}{2\pi}e^{-x^2/2}\sqrt{2\pi}
=\frac{1}{\sqrt{2\pi}}e^{-x^2/2}.
\]
So $X\sim N(0,1)$.

For $Y$,
\[
f_Y(y)=\int_{-\infty}^{\infty}\frac{1}{2\pi}
e^{-x^2/2-(x-y)^2/2}\,dx.
\]
Complete the square:
\[
\frac{x^2}{2}+\frac{(x-y)^2}{2}
=\left(x-\frac y2\right)^2+\frac{y^2}{4}.
\]
Therefore,
\[
f_Y(y)=\frac{1}{2\pi}e^{-y^2/4}
\int_{-\infty}^{\infty} e^{-(x-y/2)^2}\,dx.
\]
Since $\int_{-\infty}^{\infty}e^{-u^2}\,du=\sqrt{\pi}$,
\[
f_Y(y)=\frac{1}{2\pi}e^{-y^2/4}\sqrt{\pi}
=\frac{1}{2\sqrt{\pi}}e^{-y^2/4}
=\frac{1}{\sqrt{4\pi}}e^{-y^2/4}.
\]
Thus $Y\sim N(0,2)$.
\item If $X$ and $Y$ were independent, then $f_{X,Y}(x,y)=f_X(x)f_Y(y)$.
But
\[
f_X(x)f_Y(y)=\frac{1}{2\pi}e^{-x^2/2-y^2/4},
\]
while
\[
f_{X,Y}(x,y)=\frac{1}{2\pi}e^{-x^2/2-(x-y)^2/2}.
\]
These are not equal in general because the exponent in $f_{X,Y}$ contains the
cross term $xy$. Therefore $X$ and $Y$ are not independent.
\end{enumerate}

\finalanswer{\text{(a)}\ c=\frac{1}{2\pi}}
\finalanswer{\text{(b)}\ f_X(x)=\frac{1}{\sqrt{2\pi}}e^{-x^2/2},\quad
f_Y(y)=\frac{1}{\sqrt{4\pi}}e^{-y^2/4}}
\finalanswer{\text{(c)}\ X\ \text{and}\ Y\ \text{are not independent}}

\end{solution}


\end{document}
